% Options for packages loaded elsewhere
\PassOptionsToPackage{unicode}{hyperref}
\PassOptionsToPackage{hyphens}{url}
\documentclass[
  ignorenonframetext,
]{beamer}
\newif\ifbibliography
\usepackage{pgfpages}
\setbeamertemplate{caption}[numbered]
\setbeamertemplate{caption label separator}{: }
\setbeamercolor{caption name}{fg=normal text.fg}
\beamertemplatenavigationsymbolsempty
% remove section numbering
\setbeamertemplate{part page}{
  \centering
  \begin{beamercolorbox}[sep=16pt,center]{part title}
    \usebeamerfont{part title}\insertpart\par
  \end{beamercolorbox}
}
\setbeamertemplate{section page}{
  \centering
  \begin{beamercolorbox}[sep=12pt,center]{section title}
    \usebeamerfont{section title}\insertsection\par
  \end{beamercolorbox}
}
\setbeamertemplate{subsection page}{
  \centering
  \begin{beamercolorbox}[sep=8pt,center]{subsection title}
    \usebeamerfont{subsection title}\insertsubsection\par
  \end{beamercolorbox}
}
% Prevent slide breaks in the middle of a paragraph
\widowpenalties 1 10000
\raggedbottom
\AtBeginPart{
  \frame{\partpage}
}
\AtBeginSection{
  \ifbibliography
  \else
    \frame{\sectionpage}
  \fi
}
\AtBeginSubsection{
  \frame{\subsectionpage}
}
\usepackage{iftex}
\ifPDFTeX
  \usepackage[T1]{fontenc}
  \usepackage[utf8]{inputenc}
  \usepackage{textcomp} % provide euro and other symbols
\else % if luatex or xetex
  \usepackage{unicode-math} % this also loads fontspec
  \defaultfontfeatures{Scale=MatchLowercase}
  \defaultfontfeatures[\rmfamily]{Ligatures=TeX,Scale=1}
\fi
\usepackage{lmodern}
\ifPDFTeX\else
  % xetex/luatex font selection
\fi
% Use upquote if available, for straight quotes in verbatim environments
\IfFileExists{upquote.sty}{\usepackage{upquote}}{}
\IfFileExists{microtype.sty}{% use microtype if available
  \usepackage[]{microtype}
  \UseMicrotypeSet[protrusion]{basicmath} % disable protrusion for tt fonts
}{}
\makeatletter
\@ifundefined{KOMAClassName}{% if non-KOMA class
  \IfFileExists{parskip.sty}{%
    \usepackage{parskip}
  }{% else
    \setlength{\parindent}{0pt}
    \setlength{\parskip}{6pt plus 2pt minus 1pt}}
}{% if KOMA class
  \KOMAoptions{parskip=half}}
\makeatother
\setlength{\emergencystretch}{3em} % prevent overfull lines
\providecommand{\tightlist}{%
  \setlength{\itemsep}{0pt}\setlength{\parskip}{0pt}}
% Auto-generated by infrastructure/rendering/_pdf_latex_helpers.py::extract_math_font_preamble.
% Minimal header passed to Pandoc via -H so Beamer slides pick up the same math font
% as the combined-PDF path. Adjust the manuscript's preamble.md to change the font globally.
\usepackage{unicode-math}
\setmathfont{latinmodern-math.otf}

% Slide layout defaults for warning-clean scientific decks.
\usepackage{etoolbox}
\IfFileExists{xurl.sty}{\usepackage{xurl}}{}
\IfFileExists{seqsplit.sty}{\usepackage{seqsplit}}{\newcommand{\seqsplit}[1]{#1}}
\protected\def\breakseq#1{\seqsplit{#1}}
\protected\def\breaktt#1{\begingroup\ttfamily\seqsplit{#1}\endgroup}
\setlength{\emergencystretch}{6em}
\tolerance=5000
\hbadness=10000
\hfuzz=1pt
\setlength{\tabcolsep}{2pt}
\AtBeginEnvironment{longtable}{\tiny\renewcommand{\arraystretch}{0.86}\setlength{\tabcolsep}{1pt}}
\AtBeginEnvironment{tabular}{\tiny\renewcommand{\arraystretch}{0.86}\setlength{\tabcolsep}{1pt}}

% Natbib and cross-reference fallbacks — slides don't load natbib
% or cleveref, but combined-PDF manuscript prose may emit these
% commands. The fallback renders citations as a bracketed key list
% and cross-references as detokenized labels so slides don't fail on
% undefined control sequences or raw underscores. \providecommand is
% a no-op if packages load later.
\providecommand{\citep}[1]{[#1]}
\providecommand{\citet}[1]{#1}
\providecommand{\citealp}[1]{#1}
\providecommand{\citeauthor}[1]{#1}
\providecommand{\citeyear}[1]{#1}
\providecommand{\cref}[1]{\texttt{\detokenize{#1}}}
\providecommand{\Cref}[1]{\texttt{\detokenize{#1}}}
\usepackage{bookmark}
\IfFileExists{xurl.sty}{\usepackage{xurl}}{} % add URL line breaks if available
\urlstyle{same}
% fallback for those not using the hyperref driver hyperxmp:
\makeatletter
\@ifundefined{xmpquote}{\newcommand{\xmpquote}[1]{#1}}{}
\makeatother
\hypersetup{
  hidelinks,
  pdfcreator={LaTeX via pandoc}}

\author{\texorpdfstring{}{}}
\date{}

\begin{document}

\section{Methodology}\label{sec:methodology}

\begin{frame}[fragile,allowframebreaks]{Methodology}
The DSL is implemented as eight cooperating modules under
\breaktt{src/methods\_dsl/}, each corresponding to one stage of a
BPL-inspired pipeline \breakseq{[@bpl2026]}. This section walks the pipeline
stage by stage, naming the function or class that implements each design
decision so every claim below is directly checkable against
\breaktt{src/methods\_dsl/}.
\end{frame}

\begin{frame}[fragile,allowframebreaks]{Controlled vocabulary
(\texttt{vocabulary.py})}
\protect\phantomsection\label{controlled-vocabulary-vocabulary.py}
A \texttt{StepKind} is one of 9 controlled intents ---
\texttt{TRANSFER}, \texttt{ADD}, \texttt{MIX}, \texttt{INCUBATE},
\texttt{MEASURE}, \texttt{WAIT}, \texttt{COMPUTE}, \texttt{VALIDATE},
\texttt{ANNOTATE} --- and a \texttt{Target} is one of 3 execution
backends --- \texttt{HUMAN}, \texttt{AUTOMATED}, \texttt{SIMULATION}.
\texttt{target\_accepts} encodes which kinds require an automated
backend: only \texttt{COMPUTE} has no manual equivalent in this DSL's
scope, so \texttt{HUMAN} and \texttt{SIMULATION} accept every other
kind. This is the domain-neutral generalization of BPL's protocol-level
verbs (\texttt{add\_reagent}, \texttt{transfer}, \texttt{incubate}): a
step names \emph{what} happens, never \emph{how} a particular backend
performs it.
\end{frame}

\begin{frame}[fragile,allowframebreaks]{Dimensional safety
(\texttt{units.py})}
\protect\phantomsection\label{dimensional-safety-units.py}
Every \breaktt{Quantity(value,\ unit)} resolves its \texttt{unit} to one
of seven \texttt{Dimension} members (mass, volume, temperature, time,
concentration, count, dimensionless) via \texttt{dimension\_of}, drawing
from a controlled table of 18 unit strings. \breaktt{check\_compatible}
raises \texttt{DimensionError} the moment two quantities with different
dimensions are combined --- the concrete realization of BPL's design
principle that ``the type system catches \texttt{mL\ +\ g} at compile
time, not at the bench.'' Temperature is tracked as its own dimension
with no shared base unit (\texttt{degC} and \texttt{K} are never
auto-converted), since this DSL has no use for that conversion and an
incorrect affine conversion is worse than refusing one.
\end{frame}

\begin{frame}[fragile,allowframebreaks]{Method model
(\texttt{model.py})}
\protect\phantomsection\label{method-model-model.py}
A \texttt{Method} is a name, a version, a target, a tuple of
\texttt{Resource} declarations (anything a step reads from or writes to
--- generalizing BPL's \texttt{reagent}/\texttt{labware}), a tuple of
method-level \texttt{Parameter}s, and a tuple of \texttt{Step}s. Each
\texttt{Step} carries a \texttt{step\_id}, a \texttt{StepKind}, a
\texttt{Target}, its own parameters, an optional expected duration (must
be a time \texttt{Quantity}), and a \texttt{depends\_on} tuple of
prerequisite \texttt{step\_id}s --- the explicit DAG edges this
section's compilation stage resolves. All four dataclasses are frozen;
\texttt{\_\_post\_init\_\_} rejects malformed shapes immediately (empty
names, self-dependency, a non-time \breaktt{expected\_duration}) so
structurally invalid methods cannot even be constructed, collapsing
BPL's syntax gate into Python's own construction-time checks.
\end{frame}

\begin{frame}[fragile,allowframebreaks]{Staged validation
(\texttt{validation.py})}
\protect\phantomsection\label{staged-validation-validation.py}
\texttt{run\_all\_gates} runs exactly 4 gates in fixed order, mirroring
BPL's staged short-circuit (a syntax failure never reaches the plan
gate):

\begin{enumerate}
\tightlist
\item
  \textbf{\texttt{structural\_gate}} --- every \texttt{step\_id} is
  unique and every \texttt{depends\_on} entry resolves to a real step.
\item
  \textbf{\texttt{semantic\_gate}} --- every \texttt{Quantity} attached
  to a resource, a method parameter, a step's expected duration, or a
  step parameter resolves to a known \texttt{Dimension} (catches the
  unit-vocabulary violation a frozen dataclass's
  \texttt{\_\_post\_init\_\_} cannot, since \texttt{Quantity} does not
  validate its unit eagerly).
\item
  \textbf{\texttt{plan\_gate}} --- the step-dependency graph is acyclic,
  checked by attempting \breaktt{topological\_order} and catching
  \texttt{CycleError}.
\item
  \textbf{\texttt{target\_gate}} --- every step's target is compatible
  with the method's target (\texttt{HUMAN} methods accept only
  \texttt{HUMAN} steps; \texttt{AUTOMATED} methods accept both;
  \texttt{SIMULATION} methods accept only \texttt{SIMULATION} steps) and
  every step's kind is executable on its assigned target.
\end{enumerate}

If either of the first two gates fails, \texttt{run\_all\_gates} returns
early with only those two results --- \texttt{plan\_gate} and
\texttt{target\_gate} assume a structurally and semantically valid
method and would otherwise report misleading secondary failures.
\end{frame}

\begin{frame}[fragile,allowframebreaks]{Deterministic compilation
(\texttt{compiler.py})}
\protect\phantomsection\label{deterministic-compilation-compiler.py}
\texttt{compile\_method} first calls \texttt{run\_all\_gates} and raises
\breaktt{MethodValidationError} (carrying every failed gate's issues) if
any gate fails. On success, \breaktt{topological\_order} schedules the
validated steps with Kahn's algorithm \breakseq{[@kahn1962topological]}:
repeatedly remove a step whose dependencies are all already scheduled,
breaking ties by ascending \texttt{step\_id} so the same method always
yields the same order --- Python does not guarantee dict/set iteration
order is stable for this purpose, so the tie-break is explicit, not
incidental. The scheduled steps are then encoded as a canonical,
sort-keys JSON payload and hashed with SHA-256
(\breaktt{\_compute\_plan\_hash}) into \texttt{Plan.plan\_hash}. Because
the hash is computed purely from \texttt{method\_name},
\texttt{method\_version}, \texttt{target}, and each step's
\texttt{step\_id}/\texttt{name}/\texttt{kind}/\texttt{target}/scheduled
\texttt{order} --- never from a wall-clock timestamp or a UUID ---
recompiling the same \texttt{Method} object always produces the same
\texttt{plan\_hash}, which \breakseq{[@sec:results]} verifies live rather than
asserts.
\end{frame}

\begin{frame}[fragile,allowframebreaks]{Export (\texttt{export.py})}
\protect\phantomsection\label{export-export.py}
A compiled \texttt{Plan} renders to four formats, mirroring BPL's
``CSV/XLSX worklists, workflow graphs'' export surface at a scope
appropriate for a template exemplar: \breaktt{to\_worklist\_markdown} (a
numbered, human-readable worklist),
\texttt{to\_csv\_rows}/\texttt{write\_csv} (machine-readable rows),
\texttt{to\_mermaid} (a \texttt{flowchart\ TD} showing scheduled order),
and \texttt{to\_json}/\texttt{write\_json} (the exact canonical JSON the
plan hash was computed over, so a reader can independently verify
\texttt{Plan.plan\_hash} by re-hashing the exported file).
\end{frame}

\begin{frame}[fragile,allowframebreaks]{Provenance (\texttt{trust.py})}
\protect\phantomsection\label{provenance-trust.py}
\texttt{ProvenanceTier} orders three levels of trust ---
\texttt{DECLARED}, \texttt{CALIBRATED}, \texttt{VERIFIED} ---
generalizing BPL's audit model. \texttt{append\_record} extends an
immutable hash-chain of \texttt{StateRecord}s, each hashing its own
\texttt{key}/\texttt{value}/ \texttt{tier}/\texttt{prev\_hash}
\breakseq{[@merkle1987digital]}; \texttt{verify\_chain} recomputes every
record's hash and checks it against the recorded \texttt{prev\_hash}
chain. This is a consistency check, not a cryptographic tamper-proof
guarantee against an actor with write access to the whole stored chain:
it detects in-chain tampering (any record after a tampered one no longer
matches) but cannot detect a chain rewritten from record zero, exactly
the boundary BPL's own ``hash-chained'' (not cryptographically signed)
audit model claims.
\end{frame}

\begin{frame}[fragile,allowframebreaks]{Zero-mock testing methodology}
\protect\phantomsection\label{zero-mock-testing-methodology}
The project is governed by a strict zero-mock policy, evaluated by
running
\breaktt{uv\ run\ pytest\ projects/templates/template\_methods\_paper/tests}
during the build.

\begin{enumerate}
\tightlist
\item
  \textbf{Library tests} exercise every gate, the compiler, the
  exporters, and the trust module against real \texttt{Method} objects
  --- \texttt{conftest.py} ships one fixture per gate-failure mode
  (unknown dependency, duplicate step id, unknown unit, cyclic
  dependency, target mismatch) plus a linear-chain and a diamond-DAG
  method for scheduling tests. No \texttt{unittest.mock}, no
  \texttt{MagicMock}, no \texttt{@patch}.
\item
  \textbf{Script test} runs \breaktt{run\_methods\_analysis()} against a
  temporary output root and asserts that real worklist/CSV/Mermaid/JSON
  files, reports, and a real PNG figure are written.
\item
  \textbf{Determinism tests} compile the same \texttt{Method} object
  twice and assert \texttt{plan\_hash} equality live, rather than
  asserting against a hardcoded hash literal --- a literal would
  silently stop testing the moment the compiler's hash input changed.
\item
  \textbf{Coverage gate}: CI enforces a ≥90\% statement-coverage gate on
  \breaktt{projects/templates/template\_methods\_paper/src/}; the live
  figure is tracked in
  \href{../../../../docs/_generated/COUNTS.md}{\breaktt{docs/\_generated/COUNTS.md}}.
\end{enumerate}
\end{frame}

\end{document}
